%Daten zur Institution

%\input{Dozentdaten}

%\renewcommand{\fachbereich}{Fachbereich}

%\renewcommand{\dozent}{Prof. Dr. . }

%Klausurdaten

\renewcommand{\klausurgebiet}{ }

\renewcommand{\klausurtyp}{ }

\renewcommand{\klausurdatum}{ . 20}

\klausurvorspann {\fachbereich} {\klausurdatum} {\dozent} {\klausurgebiet} {\klausurtyp}

%Daten für folgende Punktetabelle

\renewcommand{\aeins}{  3  }

\renewcommand{\azwei}{  3   }

\renewcommand{\adrei}{  0  }

\renewcommand{\avier}{  4  }

\renewcommand{\afuenf}{  5  }

\renewcommand{\asechs}{  4  }

\renewcommand{\asieben}{  9  }

\renewcommand{\aacht}{  0  }

\renewcommand{\aneun}{  8  }

\renewcommand{\azehn}{  6  }

\renewcommand{\aelf}{  0  }

\renewcommand{\azwoelf}{  3  }

\renewcommand{\adreizehn}{  3   }

\renewcommand{\avierzehn}{  0  }

\renewcommand{\afuenfzehn}{  0  }

\renewcommand{\asechzehn}{  7  }

\renewcommand{\asiebzehn}{  55  }

\renewcommand{\aachtzehn}{    }

\renewcommand{\aneunzehn}{    }

\renewcommand{\azwanzig}{    }

\renewcommand{\aeinundzwanzig}{    }

\renewcommand{\azweiundzwanzig}{    }

\renewcommand{\adreiundzwanzig}{    }

\renewcommand{\avierundzwanzig}{    }

\renewcommand{\afuenfundzwanzig}{    }

\renewcommand{\asechsundzwanzig}{    }

\punktetabellesechzehn

\klausurnote

\newpage

\setcounter{section}{K}

\inputaufgabegibtloesung
{3}
{

Definiere die folgenden
\zusatzklammer {kursiv gedruckten} {} {} Begriffe.
\aufzaehlungsechs{Die
\stichwort {Gaußkrümmung} {}
zu einer differenzierbaren Fläche

\mavergleichskette
{\vergleichskette
{ Y
}
{ \subseteq }{ \R^3
}
{  }{
}
{    }{
}
{   }{
}
}
{}{}{}
in einem Punkt

\mavergleichskette
{\vergleichskette
{ P
}
{ \in }{ Y
}
{  }{
}
{    }{
}
{   }{
}
}
{}{}{.}

}{Eine
\stichwort {topologische Mannigfaltigkeit} {}
der Dimension $n$.

}{Ein
\stichwort {zeitunabhängiges Vektorfeld} {}
auf einer differenzierbaren Mannigfaltigkeit $M$.

}{Eine
\stichwort {Differentialform vom Grad $k$ } {}
auf einer differenzierbaren Mannigfaltigkeit $M$.

}{Eine \stichwort {riemannsche Mannigfaltigkeit} {.}

}{Ein
\stichwort {torsionsfreier} {}
linearer Zusammenhang auf dem Tangentialbündel einer riemannschen Mannigfaltigkeit $M$.
}

}
{} {}

\inputaufgabegibtloesung
{3}
{

Formuliere die folgenden Sätze.
\aufzaehlungdrei{Der
\stichwort {Existenzsatz für den Paralleltransport auf einer Hyperfläche} {.}}{Der Satz über die lokale Beschreibung von Differentialformen.}{Die
\stichwort {Quaderversion} {}
des Satzes von Stokes.}

}
{} {}

\inputaufgabe
{0}
{

}
{} {}

\inputaufgabegibtloesung
{4 (1+3)}
{

Es sei $M$ die durch

\mavergleichskettedisp
{\vergleichskette
{ x^2+3y^2
}
{ =} { 2
}
{ } {
}
{ } {
}
{ } {
}
}
{}{}{}
gegebene Ellipse im $\R^2$,

\mavergleichskette
{\vergleichskette
{ P
}
{ = }{ \begin{pmatrix}  1  \\ { \frac{   1  }{  \sqrt{3}  } }    \end{pmatrix}
}
{  }{
}
{    }{
}
{   }{
}
}
{}{}{}
und

\mavergleichskette
{\vergleichskette
{ Q
}
{ = }{ \begin{pmatrix}  -1 \\ { \frac{   1  }{  \sqrt{3}  } }    \end{pmatrix}
}
{  }{
}
{    }{
}
{   }{
}
}
{}{}{}
Punkte auf $M$.
\aufzaehlungdrei{Zeige, dass

\mavergleichskettedisp
{\vergleichskette
{ v
}
{ =} { \begin{pmatrix}  -  \sqrt{3}  \\ 1    \end{pmatrix}
}
{ } {
}
{ } {
}
{ } {
}
}
{}{}{}
ein
\definitionsverweis {Tangentialvektor}{}{}
in $P$ an $M$  ist.
}{Bestimme das Bild von $v$ unter einem Paralleltransport auf $M$ von $P$ nach $Q$.
}{
}

}
{} {}

\inputaufgabegibtloesung
{5}
{

Beweise den Satz über die Orientierungen auf einer differenzierbaren Hyperfläche.

}
{} {}

\inputaufgabegibtloesung
{4}
{

Es sei

\mavergleichskette
{\vergleichskette
{ Y
}
{ \subseteq }{ W
}
{ \subseteq }{ \R^3
}
{    }{
}
{   }{
}
}
{}{}{,}
$W$
\definitionsverweis {offen}{}{,}
eine
\definitionsverweis {differenzierbare Fläche}{}{}
versehen mit einem
\definitionsverweis {Einheitsnormalenfeld}{}{}
$N$. Es sei
\maabb {\gamma} { I } { Y
} {}
eine stetig differenzierbare Kurve und sei
\maabbdisp {F} { I } { \R^3
} {}
ein differenzierbares
\definitionsverweis {tangentiales Vektorfeld}{}{}
längs $\gamma$, das
\definitionsverweis {parallel}{}{}
sei. Zeige, dass dann auch das Vektorfeld
\maabbeledisp {} { I } { \R^3
} { t } { N(\gamma(t)) \times F(t)
} {,}
parallel ist, wobei $\times$ das
\definitionsverweis {Kreuzprodukt}{}{}
im $\R^3$ bezeichnet.

}
{} {}

\inputaufgabegibtloesung
{9}
{

Es sei

\mavergleichskette
{\vergleichskette
{ M
}
{ \neq }{ \emptyset
}
{  }{
}
{    }{
}
{   }{
}
}
{}{}{}
eine
\definitionsverweis {differenzierbare Mannigfaltigkeit}{}{}
der Dimension $n$. Zeige, dass es eine Kette von
\definitionsverweis {abgeschlossenen Untermannigfaltigkeiten}{}{}

\mavergleichskettedisp
{\vergleichskette
{ M_0
}
{ \subseteq} { M_1
}
{ \subseteq} { M_2
}
{ \subseteq  \ldots \subseteq} { M_{n-1}
}
{ \subseteq} { M_n
}
}
{
\vergleichskettefortsetzung
{ =} { M
}
{ } {}
{ } {}
{ } {}
}{}{}
derart gibt, dass die abgeschlossene Untermannigfaltigkeit $M_i$ die Dimension $i$ besitzt.

}
{} {}

\inputaufgabe
{0}
{

}
{} {}

\inputaufgabegibtloesung
{8}
{

Es sei
\maabbdisp {\psi} {\R^n } {\R
} {}
eine
\definitionsverweis {differenzierbare Funktion}{}{}
und

\mavergleichskettedisphandlinks
{\vergleichskettedisphandlinks
{M
}
{ =} { \left\{  \left(  x_1   , \,   , \ldots ,  , \,   x_n , \,   \psi(x_1 , \ldots , x_n)                \right)   \mid  (x_1 , \ldots , x_n) \in \R^n         \right\}
}
{ \subseteq} { \R^n \times \R
}
{ } {
}
{ } {
}
}
{}{}{}
der
\definitionsverweis {Graph}{}{}
von $\psi$. Zeige, dass $M$ eine
\definitionsverweis {orientierte}{}{} \definitionsverweis {riemannsche Mannigfaltigkeit}{}{}
ist und dass für die
\definitionsverweis {kanonische Volumenform}{}{}
$\omega$ auf $M$ die Formel

\mavergleichskettedisp
{\vergleichskette
{ (\operatorname{Id} \times \psi)^*(\omega)
}
{ =} { \left(  1+ \sum_{i = 1}^n \left(  { \frac{   \partial   \psi  }{  \partial   x_i   } }  \right)^2  \right)^{1/2} dx_1 \wedge \ldots \wedge dx_n
}
{ } {
}
{ } {
}
{ } {
}
}
{}{}{}
gilt.

}
{} {}

\inputaufgabe
{6}
{

Berechne die Oberfläche der Einheitskugel.

}
{} {}

\inputaufgabe
{0}
{

}
{} {}

\inputaufgabegibtloesung
{3}
{

Berechne die
\definitionsverweis {äußere Ableitung}{}{}
$d \omega$ der
$1$-\definitionsverweis {Differentialform}{}{}

\mavergleichskettedisp
{\vergleichskette
{ \omega
}
{ =} { { \frac{   x^2 }{  y } } dx- { \frac{   x  }{  y^2 } } dy
}
{ } {
}
{ } {
}
{ } {
}
}
{}{}{}
auf

\mavergleichskette
{\vergleichskette
{ U
}
{ = }{ \left\{ (x,y) \in \R^2  \mid   y \neq 0         \right\}
}
{  }{
}
{    }{
}
{   }{
}
}
{}{}{.}

}
{} {}

\inputaufgabegibtloesung
{3}
{

Man gebe ein Beispiel für eine
\definitionsverweis {zusammenhängende}{}{}
\definitionsverweis {Mannigfaltigkeit mit Rand}{}{,}
deren Rand aus
\zusatzklammer {einer disjunkten Vereinigung von} {} {}
unendlich vielen Kreisen $S^1$ besteht.

}
{} {}

\inputaufgabe
{0}
{

}
{} {}

\inputaufgabe
{0}
{

}
{} {}

\inputaufgabegibtloesung
{7}
{

Es sei $M$ eine
\definitionsverweis {riemannsche Mannigfaltigkeit}{}{}
und sei ein
\definitionsverweis {linearer Zusammenhang}{}{}
auf dem
\definitionsverweis {Tangentialbündel}{}{}
$TM$ gegeben, der
\definitionsverweis {metrisch}{}{}
und
\definitionsverweis {torsionsfrei}{}{}
sei. Zeige, dass für Vektorfelder $X,Y,Z$ die Formel

\mavergleichskettedisp
{\vergleichskette
{ \left\langle \nabla_X Y , Z  \right\rangle
}
{ =} { { \frac{   1  }{  2 } }  \left(  D_X \left\langle  Y  ,  Z  \right\rangle +  D_Y \left\langle  Z  ,  X  \right\rangle -  D_Z \left\langle  X  ,  Y  \right\rangle  - \left\langle  Y  , [X,Z]   \right\rangle - \left\langle  Z  , [Y,X]   \right\rangle + \left\langle  X  , [Z,Y]   \right\rangle   \right)
}
{ } {
}
{ } {
}
{ } {
}
}
{}{}{}
gilt.

}
{} {}

 [[Kategorie:Latexseite]]
